\section*{Hoofdstuk \ref{ch:b1:ocean-circulation} --- De circulatie van de oceaan}
% values: calc/ch29.py

\begin{solution}{exo:b1:ocean-circulation:1}
(a)~Bij \qty{4}{m/s} (eerste formule): $1000\,C_D = 0.29 + 0.775 + 0.481 = 1.546$,
$T = 1.3 \times \num{1.546e-3} \times 16 = \qty{0.0322}{Pa}$. Bij \qty{7}{m/s}
(tweede formule): $1000\,C_D = 0.60 + 0.497 = 1.097$, $T = \qty{0.0699}{Pa}$. Bij
\qty{15}{m/s}: $1000\,C_D = 1.665$, $T = \qty{0.487}{Pa}$. (b)~$0.487/0.0699 =
6.97$, tegen $(15/7)^2 = 4.59$: de weerstandscoëfficiënt zelf groeit met de wind,
omdat een sterkere wind ruwere golven opwekt die de lucht beter vasthouden.
\end{solution}

\begin{solution}{exo:b1:ocean-circulation:2}
(a)~Op het noordelijk halfrond, \ang{90} naar rechts van een wind die naar het oosten waait: naar het zuiden. (b)~De wind waait naar het zuidwesten; rechts ervan ligt het noordwesten. (c)~Op het zuidelijk halfrond, \ang{90} naar links van een wind die naar het oosten waait: naar het noorden. (d)~De wind waait naar het noorden; links ervan ligt het westen, dat op een westkust naar zee wijst. In geval (d) wijst het transport naar zee: het oppervlaktewater wordt van onderaf vervangen, en koud water stijgt langs de kust op.
\end{solution}

\begin{solution}{exo:b1:ocean-circulation:3}
(a)~$31.7 + 3.5 - 17.9 = \qty{17.3}{\sverdrup}$, tegen de gepubliceerde
\qty{17.1}{\sverdrup}: de gemiddelden van de drie delen tellen binnen de afronding en de onzekerheden op tot de gemiddelde omwenteling ($\pm\qty{3.2}{\sverdrup}$).
(b)~Boven \qty{1100}{m}: $16.6 + 0.4 = \qty{17.0}{\sverdrup}$ noordwaarts; eronder: $12.0 + 5.9 = \qty{17.9}{\sverdrup}$ zuidwaarts. De twee heffen elkaar bijna op, zoals moet wanneer er geen water zich ten noorden van de doorsnede ophoopt; het kleine verschil ligt binnen de onzekerheden. (c)~$\num{17.1e6} \times \num{3.156e7} = \qty{5.40e14}{m^3}$, dat wil zeggen
\qty{5.40e5}{km^3} per jaar; $\num{1.3324e9}/(\num{5.40e5}) \approx 2470$ jaar.
\end{solution}

\begin{solution}{exo:b1:ocean-circulation:4}
$D_E = 7.6 \times 10/\sqrt{\sin\varphi}$ en $V_0 = 0.0127 \times
10/\sqrt{\sin\varphi}$. Bij \ang{15}: \qty{149}{m} en \qty{0.250}{m/s}
(\qty{2.5}{\percent} van de wind); bij \ang{45}: \qty{90}{m} en \qty{0.151}{m/s}
(\qty{1.5}{\percent}); bij \ang{60}: \qty{82}{m} en \qty{0.136}{m/s}
(\qty{1.4}{\percent}). Nabij de evenaar gaat $\sin\varphi \to 0$ en divergeren beide formules; fundamenteler: $f \to 0$, zodat de Coriolis-kracht die nodig is om de wrijving in evenwicht te houden verdwijnt, en de tijd $1/f$ die het evenwicht nodig heeft om zich te vestigen langer wordt dan de wind stabiel blijft. Ekman beperkte zijn formule tot $|\varphi| \geq
\ang{10}$.
\end{solution}

\begin{solution}{exo:b1:ocean-circulation:5}
(a)~$f = 2\Omega\sin\ang{11} = \qty{2.783e-5}{s^{-1}}$; $M = 0.07/(\num{2.783e-5}) =
\qty{2515}{kg.m^{-1}.s^{-1}}$, \ang{90} naar rechts van een wind die naar het westen waait:
naar het noorden. (b)~$Q = (2515/1027) \times \num{5e6} = \qty{1.22e7}{m^3/s}$,
\qty{12.2}{\sverdrup}. (c)~De schatting ligt binnen de onzekerheid van de rechtstreekse
meting ($12.0 \pm 5.5$) en dicht bij de waarde uit langjarige winden
($13.5 \pm 0.3$): de stelling verklaart het waargenomen transport.
\end{solution}

\begin{solution}{exo:b1:ocean-circulation:6}
(a)~$f(\ang{2}) = \qty{5.09e-6}{s^{-1}}$; $M = 0.05/(\num{5.09e-6}) =
\qty{9824}{kg.m^{-1}.s^{-1}}$, \qty{9.57}{m^2/s}: op \ang{2} N naar het noorden
(rechts van een westenwind), op \ang{2} Z naar het zuiden (links). (b)~Over een strook van
\qty{1}{m} lang in de richting van de lengtegraden, $2 \times 9.57 = \qty{19.1}{m^3/s}$ vertrekken over de twee
breedtegraden, \qty{444.8}{km} uit elkaar: $w = 19.1/(\num{4.448e5}) = \qty{4.30e-5}{m/s}$,
ongeveer \qty{3.7}{m} per dag. (c)~Op de evenaar is $f = 0$ en is $|\vec M| = T/|f|$
oneindig: het stationaire Ekmanevenwicht bestaat daar niet, en de getallen hierboven, zeer gevoelig voor de gekozen \omterm{def:b1:spherical-astronomy:geographic}{breedtegraad}, zijn slechts een orde van grootte. Wat overeind blijft is het teken: aan beide zijden wijst het transport weg van de evenaar, zodat het oppervlaktewater divergeert en koud water langs de evenaar opstijgt.
\end{solution}

\begin{solution}{exo:b1:ocean-circulation:7}
(a)~$f(\ang{15}) = \qty{3.775e-5}{s^{-1}}$: $M = 0.06/(\num{3.775e-5}) =
\qty{1590}{kg.m^{-1}.s^{-1}}$ naar het noorden; $f(\ang{40}) =
\qty{9.375e-5}{s^{-1}}$: $M = \qty{1280}{kg.m^{-1}.s^{-1}}$ naar het zuiden.
(b)~Tussen de twee breedtegraden, \qty{2780}{km} uit elkaar, $\partial M_y/\partial y =
-(1590 + 1280)/(\num{2.78e6}) = \qty{-1.03e-3}{kg.m^{-2}.s^{-1}}$: volgens
\cref{prop:b1:ocean-circulation:pumping} is $w_E = \qty{-1.005e-6}{m/s}$, een neerwaartse
snelheid van \qty{31.7}{m} per jaar. (c)~$\num{1.005e-6} \times \num{2.78e6} \times
\num{5e6} = \qty{1.40e7}{m^3/s}$, \qty{14.0}{\sverdrup} naar beneden gepompt.
\end{solution}

\begin{solution}{exo:b1:ocean-circulation:8}
(a)~$H = 1027 \times 3992 \times \num{17e6} \times 12 = \qty{8.36e14}{W}$,
\qty{0.84}{PW}. (b)~Dat is \qty{70}{\percent} van de gemeten \qty{1.2}{PW}: de werkelijke bovenste tak is nabij het oppervlak warmer dan het gemiddelde van \qty{15}{\celsius}, en de \omterm{def:b1:ocean-circulation:gyre}{gyre} voert ook warmte. (c)~$1027 \times 3992 \times \num{17e6} \times 3 =
\qty{2.1e14}{W}$, \qty{0.21}{PW}, een kwart van (a). Alleen water dat koud wordt voordat het terugkeert, draagt veel warmte: het warmtetransport van de Atlantische Oceaan hangt af van de \omterm{def:b1:ocean-circulation:moc}{vorming van diep water}.
\end{solution}

\begin{solution}{exo:b1:ocean-circulation:9}
Laat $x$, $y$ en $1 - x - y$ de aandelen van AAIW, NADW en AABW zijn. Zout- en warmtebalansen, gemeten vanaf AABW:
\[
  -0.52\,x + 0.33\,y = 34.731 - 34.67 = 0.061,\qquad
  2.9\,x + 3.6\,y = 2.78 - 0.4 = 2.38 .
\]
De determinant is $-0.52 \times 3.6 - 0.33 \times 2.9 = -2.829$, en de regel van Cramer
geeft $x = (0.061 \times 3.6 - 0.33 \times 2.38)/(-2.829) = 0.20$ en $y = (-0.52
\times 2.38 - 2.9 \times 0.061)/(-2.829) = 0.50$: het monster is \qty{20}{\percent}
AAIW, \qty{50}{\percent} NADW en \qty{30}{\percent} AABW, een punt binnen de driehoek van de drie wateren.
\end{solution}

\begin{solution}{exo:b1:ocean-circulation:10}
(a)~Met $W = u + iv = V_0e^{i\pi/4}e^{(1+i)az}$, waarvan het reële en imaginaire deel de
formules van de stelling zijn, $W'' = (1+i)^2a^2W = 2ia^2W$. De twee vergelijkingen
combineren tot $A_z(u'' + iv'') = -fv + ifu = if(u + iv)$, dat wil zeggen $A_zW'' = ifW$,
wat geldt omdat $2ia^2A_z = if$ voor $a^2 = f/(2A_z)$. (b)~$W'(0) =
(1+i)aV_0e^{i\pi/4} = \sqrt2\,aV_0\,e^{i\pi/2} = i\sqrt2\,aV_0$: de spanning aan het oppervlak
$\rho A_zW'(0)$ wijst naar het noorden, met de grootte $\rho A_z\sqrt2\,aV_0 = \rho A_z\sqrt2
\sqrt{f/(2A_z)}\;T/(\rho\sqrt{fA_z}) = T$. (c)~$\int_{-\infty}^0 W\,\mathrm dz =
V_0e^{i\pi/4}/((1+i)a) = V_0/(\sqrt2\,a)$, een reëel getal: het transport heeft geen
noordelijk deel, en zijn oostelijk deel is $\rho V_0/(\sqrt2\,a) =
T/\sqrt{fA_z}\cdot\sqrt{A_z/f} = T/f$. Dit is de stelling: grootte $T/f$, \ang{90} naar rechts van een wind die naar het noorden waait.
\end{solution}

\begin{solution}{exo:b1:ocean-circulation:11}
(a)~Een noordwaartse stap $\mathrm dy = a_\oplus\,\mathrm d\varphi$ verandert $f = 2\Omega
\sin\varphi$ met $2\Omega\cos\varphi\,\mathrm d\varphi$, zodat $\beta = 2\Omega
\cos\varphi/a_\oplus = 2 \times \num{7.292e-5} \times \cos\ang{30}/(\num{6.371e6}) =
\qty{1.98e-11}{m^{-1}.s^{-1}}$. (b)~$R/\beta = \num{e-6}/(\num{1.98e-11}) =
\qty{5.04e4}{m}$, ongeveer \qty{50}{km}. (c)~Inwendig: $\num{30e6}/(\num{5e6} \times
1000) = \qty{6.0e-3}{m/s}$, \qty{6}{mm/s}; randstroom: $\num{30e6}/(\num{5.04e4}
\times 1000) = \qty{0.60}{m/s}$, honderd keer sneller, de verhouding van de breedtes. De
Golfstroom, met \qty{1.8}{m/s}, is nog sneller: de werkelijke stroom is smaller dan
deze schatting, haar stroming is nabij het oppervlak geconcentreerd, en haar terugstroming draagt bij aan haar transport.
\end{solution}

\begin{solution}{exo:b1:ocean-circulation:12}
(a)~Waar de zee dieper is dan \qty{1000}{m}, bevat de kolom boven \qty{1000}{m} hoogstens \qty{1}{km} water, en elders minder: het volume boven \qty{1000}{m} is hoogstens $\num{361.84e6} \times 1 = \qty{3.618e8}{km^3}$, zodat het volume eronder ten minste $\num{1.3324e9} - \num{3.618e8} = \qty{9.71e8}{km^3}$ is. (b)~$\num{26e6} \times \num{3.156e7} = \qty{8.21e14}{m^3}$ per jaar,
\qty{8.21e5}{km^3}; $\num{9.71e8}/(\num{8.21e5}) \approx 1200$ jaar, dezelfde orde als de duizend jaar van het transportband. (c)~$(1/2)^{25/12.3} = 0.244$, ongeveer een kwart; $\num{e-3} \times 25 \times \num{3.156e7} = \qty{7.9e5}{m}$, ongeveer \qty{790}{km}. Na 25 jaar is het tritium niet verder gekomen dan een trage drift het zou dragen, en toch is een kwart ervan nog te meten.
\end{solution}

\begin{solution}{pb:b1:ocean-circulation:1}
\begin{enumerate}
  \item Een wind uit het noorden waait naar het zuiden, langs de kust.
  \item $1000\,C_D = 0.60 + 0.071 \times 8 = 1.168$: $C_D = \num{1.168e-3}$.
  \item $T = 1.3 \times \num{1.168e-3} \times 64 = \qty{0.0972}{Pa}$.
  \item $C_D = \num{1.452e-3}$ en $T = \qty{0.272}{Pa}$: de spanning groeit met $2.80$,
        meer dan $(12/8)^2 = 2.25$, omdat $C_D$ zelf met de wind toeneemt.
  \item $\Omega = \qty{7.292e-5}{rad.s^{-1}}$; $f = 2\Omega\sin\ang{38} =
        \qty{8.979e-5}{s^{-1}}$.
  \item $2\pi/f = \qty{7.00e4}{s}$, \qty{19.4}{h}. De Coriolis-kracht draait een stroming in deze tijd één volle cirkel; voordat de wind zo lang heeft geblazen, slingert het water in plaats van zich te vestigen in het stationaire evenwicht van de stelling.
  \item \ang{90} naar rechts van een wind die naar het zuiden waait: naar het westen, de open zee in.
  \item $M = T/f = 0.0972/(\num{8.979e-5}) = \qty{1082}{kg.m^{-1}.s^{-1}}$.
  \item $1082/1027 = \qty{1.054}{m^2/s}$ per meter kust.
  \item $D_E = 7.6 \times 8/\sqrt{\sin\ang{38}} = \qty{77.5}{m}$.
  \item $1.054/77.5 = \qty{0.0136}{m/s}$, \qty{1.36}{cm/s}.
  \item $V_0 = 0.0127 \times 8/\sqrt{\sin\ang{38}} = \qty{0.129}{m/s}$, \ang{45} naar rechts van de wind: naar het zuidwesten. De drijvers geven ongeveer de helft van de waarde van Ekman, enkele \qty{0.07}{m/s}.
  \item Het water dat in de \omterm{def:b1:ocean-circulation:ekman}{Ekmanlaag} van de kust wegstroomt, kan niet over de kust worden aangevuld, en ook niet van bovenaf: het wordt vervangen door water dat nabij de kust van onder de laag omhoog stijgt,
        zoals in de doorsnede die getekend is in \cref{sec:b1:ocean-circulation:upwelling}; de isothermen buigen omhoog naar de kust toe.
  \item $1.054 \times 1000 \times 86400 = \qty{9.11e7}{m^3}$ per dag per kilometer kust.
  \item $w = 1.054/(\num{2e4}) = \qty{5.27e-5}{m/s}$, \qty{4.55}{m} per dag.
  \item $77.5/4.55 = 17$ dagen.
  \item De lucht wordt van onderaf afgekoeld door het koude water: de laag nabij het oppervlak wordt koud en stabiel, de vochtigheid kan condenseren tot mist en lage stratus, en zonder convectie is er weinig regen aan de kust.
  \item Het opgewelde water is rijk aan voedingsstoffen; zij voeden het fytoplankton, dat het zooplankton en de vis voedt.
  \item Een wind die naar het noorden waait drijft een transport \ang{90} naar rechts, naar het oosten: oppervlaktewater hoopt zich op tegen de kust en zinkt (\omterm{def:b1:ocean-circulation:upwelling}{neerwelling}); de \omterm{def:b1:ocean-circulation:upwelling}{opwelling} stopt en warmer oppervlaktewater keert naar de kust terug.
  \item $f(\ang{20}) = \qty{4.988e-5}{s^{-1}}$: $M = \qty{1948}{kg.m^{-1}.s^{-1}}$,
        $\sin\ang{38}/\sin\ang{20} = 1.80$ keer meer.
  \item Op het zuidelijk halfrond ligt het transport \ang{90} naar links van de wind. Om het naar het westen, naar zee toe, te duwen, moet de wind naar het noorden waaien, uit het zuiden: in beide gevallen waait de wind langs de kust naar de evenaar.
  \item $1000\,C_D = 0.29 + 0.62 + 0.308 = 1.218$; $T = \qty{0.0396}{Pa}$; $M =
        \qty{441}{kg.m^{-1}.s^{-1}}$; de verhouding is $0.407$ voor beide.
  \item Per meter kust: $M c_p\,\Delta T = 1082 \times 3992 \times 6 =
        \qty{2.59e7}{W}$, \qty{25.9}{MW}; langs \qty{1000}{km}, \qty{2.59e13}{W},
        \qty{25.9}{TW}.
  \item $\num{31.7e6}/1.054 = \qty{3.01e7}{m}$, ongeveer \qty{30000}{km}: \omterm{def:b1:ocean-circulation:upwelling}{kustopwelling} verplaatst veel minder water dan een \omterm{def:b1:ocean-circulation:wbc}{westelijke randstroom}.
  \item $Q = 1.054 \times \num{e6} = \qty{1.05e6}{m^3/s}$: de \omterm{def:b1:ocean-circulation:upwelling}{opwelling} langs \qty{1000}{km} kust is \qty{1.05}{\sverdrup}.
\end{enumerate}
\end{solution}
